\section{Технический проект}

\subsection{Архитектура программной системы}

Разрабатываемая информационно-аналитическая система построена по архитектуре «клиент-сервер» с использованием REST-подхода.

Серверная часть реализована на языке Python с использованием фреймворка FastAPI и обеспечивает обработку запросов, выполнение бизнес-логики, а также взаимодействие с базой данных PostgreSQL и файловым хранилищем документов.

Клиентская часть представляет собой веб-приложение, доступ к которому осуществляется через браузер. Взаимодействие между клиентом и сервером выполняется по протоколу HTTP с использованием формата JSON для обмена данными.

Архитектура системы включает следующие основные компоненты:

\begin{enumerate}
	\item Клиентское веб-приложение (интерфейс пользователя).
	\item Серверное приложение (FastAPI).
	\item Система управления базами данных PostgreSQL.
	\item Файловое хранилище для загружаемых и сгенерированных документов.
	\item Модуль анализа и проверки документов формата DOCX.
\end{enumerate}

Логика работы системы организована следующим образом. Пользователь через веб-интерфейс выполняет действия (авторизация, загрузка документа, запуск проверки и др.), которые инициируют HTTP-запросы к серверу. Сервер обрабатывает запрос, при необходимости обращается к базе данных и файловой системе, выполняет проверку документа и возвращает результат клиенту.

Модуль проверки документов осуществляет анализ структуры и оформления файлов формата DOCX, выявляет несоответствия установленным требованиям и формирует отчёт об ошибках. При необходимости система генерирует исправленный файл с визуальным выделением нарушений.

Все загружаемые документы сохраняются в файловой системе, а информация о них (путь к файлу, результаты проверки, пользователь и др.) — в базе данных PostgreSQL. Это обеспечивает разделение данных и повышает масштабируемость системы.

Для обеспечения безопасности и управления доступом используется механизм аутентификации пользователей. После успешного входа пользователь получает токен доступа, который используется при последующих запросах к API.

Общая схема взаимодействия компонентов системы представлена на рисунке~\ref{fig:architecture}.

\begin{figure}[H]
	\centering
	\includegraphics[width=1\textwidth]{architecture.jpg}
	\caption{Общая схема архитектуры системы}
	\label{fig:architecture}
\end{figure}


\subsection{Обоснование выбора средств разработки}

При выборе технологического стека учитывались следующие критерии: возможность масштабирования системы, удобство разработки и сопровождения, поддержка работы с документами формата DOCX, наличие инструментов для построения REST API, а также обеспечение надёжного хранения и обработки данных.

\subsubsection{Язык программирования Python}

Основным языком разработки выбран Python. Данный язык обладает широким набором библиотек для работы с веб-приложениями, базами данных и обработкой документов.

Ключевым преимуществом Python является наличие библиотек для анализа и модификации файлов формата DOCX. Это позволило реализовать автоматическую проверку структуры и оформления документов, а также генерацию отчётов об ошибках.

Дополнительным преимуществом является высокая скорость разработки, читаемость кода и большое сообщество, что упрощает сопровождение и развитие системы.

\subsubsection{Фреймворк FastAPI}

В качестве серверного фреймворка выбран FastAPI — современный высокопроизводительный веб-фреймворк для создания REST API.

FastAPI обеспечивает:
\begin{itemize}
	\item автоматическую валидацию входных данных;
	\item высокую производительность за счёт использования асинхронных механизмов;
	\item простую интеграцию с другими компонентами системы.
\end{itemize}

Использование FastAPI позволило реализовать серверную часть системы в виде набора API-методов, обеспечивающих:
\begin{itemize}
	\item авторизацию пользователей;
	\item загрузку документов;
	\item запуск проверки;
	\item получение результатов анализа;
	\item управление данными.
\end{itemize}

\subsubsection{СУБД PostgreSQL}

В качестве системы управления базами данных выбрана PostgreSQL.

Данная СУБД обеспечивает:
\begin{itemize}
	\item надёжное хранение данных;
	\item поддержку транзакций;
	\item работу с большими объёмами информации;
	\item обеспечение целостности данных за счёт первичных и внешних ключей.
\end{itemize}

В базе данных хранятся:
\begin{itemize}
	\item пользователи системы;
	\item информация о загруженных документах;
	\item результаты проверок;
	\item ошибки и замечания;
	\item токены авторизации.
\end{itemize}

Использование PostgreSQL позволяет эффективно обрабатывать запросы и масштабировать систему при увеличении нагрузки.

\subsubsection{Модуль обработки документов DOCX}

Для работы с файлами формата DOCX используются специализированные библиотеки Python, позволяющие:
\begin{itemize}
	\item извлекать структуру документа;
	\item анализировать форматирование (шрифты, отступы, интервалы);
	\item проверять соответствие требованиям оформления;
	\item формировать список ошибок;
	\item модифицировать документ с выделением нарушений.
\end{itemize}

Модуль реализован как отдельный компонент серверной части и интегрирован с API. Это обеспечивает его независимость и возможность дальнейшего расширения.

\subsubsection{Файловое хранилище}

Для хранения загружаемых и сгенерированных документов используется файловая система.

Такой подход позволяет:
\begin{itemize}
	\item снизить нагрузку на базу данных;
	\item хранить большие файлы без ограничений;
	\item обеспечить быстрый доступ к документам.
\end{itemize}

В базе данных при этом сохраняются только пути к файлам и связанная метаинформация.

\subsubsection{Прочие библиотеки}

В процессе разработки используются дополнительные библиотеки:

\begin{enumerate}
	\item Библиотеки для работы с DOCX — для анализа и генерации документов.
	\item Средства хеширования паролей — для обеспечения безопасности.
	\item Модуль логирования — для фиксации ошибок и событий системы.
	\item Библиотеки для работы с JSON — для обмена данными между клиентом и сервером.
\end{enumerate}

Использование выбранного технологического стека позволило реализовать масштабируемую, надёжную и удобную в использовании систему с возможностью дальнейшего расширения функциональности.

\subsection{Клиентская часть}

Клиентская часть информационно-аналитической системы представлена веб-приложением, доступ к которому осуществляется через браузер пользователя. В отличие от классических десктопных решений, управление системой полностью реализовано через веб-интерфейс, что упрощает развёртывание и обеспечивает доступ с различных устройств.

Клиентская часть отвечает за взаимодействие пользователя с системой и включает формы ввода данных, отображение результатов проверки документов, работу с архивом и навигацию по функциональным модулям системы.

На рисунке~\ref{fig:client_components} представлена диаграмма компонентов клиентской части.

\begin{figure}[H]
	\centering
	\includegraphics[scale=0.75]{components.png}
	\caption{Диаграмма компонентов клиентской части}
	\label{fig:client_components}
\end{figure}

Клиентская часть реализована в виде одностраничного веб-приложения (SPA), структура которого включает набор статических файлов, загружаемых браузером пользователя при открытии системы.

Диаграмма компонентов отражает физическую организацию клиентской части и взаимосвязи между основными файлами приложения. В качестве основного HTML-документа используется файл index.html, который подключает таблицу стилей styles.css, основной сценарий app.js, а также графические ресурсы, расположенные в каталоге images/.

Описание файлов клиентской части приведено в таблице~\ref{tab:client_files}.

\begin{xltabular}{\textwidth}{|l|X|X|}
	\caption{Описание файлов клиентской части\label{tab:client_files}}\\
	\hline
	Файл & Назначение & Содержание \\ \hline
	\endfirsthead
	
	\continuecaption{Продолжение таблицы \ref{tab:client_files}}
	\hline
	Файл & Назначение & Содержание \\ \hline
	\endhead
	
	\hline
	\endlastfoot
	
	index.html &
	Главная страница приложения &
	Содержит разметку всех функциональных разделов системы: авторизация, профиль пользователя, проверка документов, генерация титульных листов, статистика, отправка документов в РИО и административная панель. Выполняет подключение таблицы стилей и JavaScript-модуля. \\ \hline
	
	styles.css &
	Таблица стилей интерфейса &
	Определяет оформление пользовательского интерфейса: цветовую схему, шрифты, размеры элементов, анимации, адаптивную вёрстку, стили прогресс-баров, индикаторов загрузки и модальных окон. \\ \hline
	
	app.js &
	Основной сценарий клиентской части &
	Реализует бизнес-логику веб-приложения: обработку пользовательских событий, выполнение HTTP-запросов к REST API, управление токенами авторизации, динамическое обновление содержимого страниц и обработку результатов проверки документов. \\ \hline
	
	images/ &
	Каталог графических ресурсов &
	Содержит изображения и иконки интерфейса: логотипы, графические элементы навигации, анимации загрузки и вспомогательные визуальные компоненты. \\ \hline
	
\end{xltabular}

Все перечисленные файлы располагаются в каталоге статических ресурсов приложения и передаются пользователю серверной частью FastAPI. После загрузки страницы основная логика работы системы выполняется на стороне клиента без полной перезагрузки интерфейса.



\subsection{Разработка клиентской части системы}

Клиентская часть информационно-аналитической системы редакционно-издательского отдела реализована в виде веб-интерфейса, доступного пользователям через современные браузеры (Google Chrome, Mozilla Firefox, Microsoft Edge). Такой подход обеспечивает кроссплатформенность решения и исключает необходимость установки дополнительного программного обеспечения на стороне пользователя. Веб-интерфейс построен по принципу одностраничного приложения (SPA), что позволяет динамически обновлять содержимое без перезагрузки страницы, ускоряя взаимодействие и снижая нагрузку на сервер. 


Все ключевые операции — авторизация, загрузка и проверка документов, генерация титульных листов, формирование статистики, отправка материалов в РИО — вынесены в отдельные функциональные блоки, доступ к которым организован через единое навигационное меню. Интерфейс адаптирован для работы на устройствах с разным разрешением экрана благодаря использованию медиа-запросов и гибкой вёрстке на основе CSS Grid и Flexbox.

\subsubsection{Архитектура клиентской части}

Клиентская часть построена по принципу одностраничного приложения (Single Page Application, SPA). Взаимодействие с сервером осуществляется посредством асинхронных HTTP-запросов, что позволяет динамически обновлять содержимое страницы без её полной перезагрузки. Это повышает отзывчивость интерфейса и сокращает время ожидания пользователя.

Основными компонентами клиентской части являются:

\begin{enumerate}
	\item Интерфейс аутентификации — формы входа и регистрации пользователей.
	\item Интерфейс профиля - для зарегистрированных пользователей. 
	\item Интерфейс проверки документов — форма загрузки файлов, панель управления проверкой и область отображения результатов.
	\item Интерфейс генерации титульных листов — динамические формы с возможностью добавления нескольких авторов, рецензентов и направлений подготовки.
	\item Интерфейс статистики — формы для задания параметров отчётности и визуализация результатов.
	\item Интерфейс отправки документов в РИО — пошаговый мастер загрузки файлов и отправки на электронную почту.
	\item Административная панель — инструменты для создания факультетов и кафедр, доступные только пользователям с ролью администратора.
\end{enumerate}

Для отображения структуры пользовательского интерфейса и переходов между основными страницами веб-приложения была разработана диаграмма навигации. Диаграмма отражает доступные пользователю разделы системы и взаимосвязи между ними.

На рисунке~\ref{fig:navigation_diagram} схема взаимодействия маршрутизаторов серверной части.

\begin{figure}[H]
	\centering
	\includegraphics[width=1\textwidth]{navigation_diagram.png}
	\caption{Схема взаимодействия маршрутизаторов серверной части}
	\label{fig:navigation_diagram}
\end{figure}

\subsubsection{Реализация пользовательского интерфейса}

Для реализации веб-интерфейса использованы стандартные веб-технологии:

\begin{enumerate}
	\item HTML5 — для разметки страниц и структурирования содержимого.
	\item CSS3 — для стилизации элементов и обеспечения адаптивного дизайна.
	\item JavaScript (ES6+) — для реализации интерактивности и взаимодействия с сервером.
\end{enumerate}

Основные элементы пользовательского интерфейса представлены в таблице~\ref{tab:ui_components}. В таблице приведены ключевые компоненты интерфейса и их назначение в системе.

\begin{xltabular}{\textwidth}{|c|l|l|X|}
	\caption{Описание интерфейсных компонентов клиентской части\label{tab:ui_components}}\\
	\hline
	№ & Тип компонента & Имя компонента & Назначение \\ \hline
	\endfirsthead
	
	\continuecaption{Продолжение таблицы \ref{tab:ui_components}}
	\hline
	№ & Тип компонента & Имя компонента & Назначение \\ \hline
	\endhead
	
	\hline
	\endfoot
	
	1 & Форма & login\_form & Авторизация пользователя в системе \\ \hline
	2 & Форма & register\_form & Регистрация нового пользователя \\ \hline
	3 & Кнопка & upload\_btn & Загрузка документа на проверку \\ \hline
	4 & Панель & check\_subnav\_card & Переключение между типами документов (методические указания, учебные пособия, монографии) \\ \hline
	5 & Область & drop\_area & Перетаскивание файла для загрузки (drag-and-drop) \\ \hline
	6 & Индикатор & check\_loading & Отображение процесса проверки (анимация загрузки) \\ \hline
	7 & Прогресс-бар & progress\_bar & Отображение текущей страницы проверки документа \\ \hline
	8 & Таблица & result\_table & Отображение результатов проверки документа \\ \hline
	9 & Кнопка & download\_report\_btn & Скачивание PDF-отчёта о проверке \\ \hline
	10 & Кнопка & download\_docx\_btn & Скачивание документа с выделенными ошибками \\ \hline
	11 & Форма & title\_generator\_form & Генерация титульного листа (с выбором типа документа) \\ \hline
	12 & Панель & admin\_panel & Управление факультетами и кафедрами \\ \hline
	13 & Форма & statistics\_form & Формирование статистических отчётов \\ \hline
	14 & Индикатор & loading\_spinner & Отображение процесса обработки запроса \\ \hline
\end{xltabular}

Все стилевые решения оформлены в виде централизованной таблицы стилей, что обеспечивает единообразие внешнего вида всех страниц приложения. В таблице стилей определены цветовые переменные, типографика, отступы и анимационные эффекты. Применена методология адаптивной вёрстки с использованием медиа-запросов, благодаря чему интерфейс корректно отображается на устройствах с различной шириной экрана (от мобильных телефонов до широкоформатных мониторов).

Особое внимание уделено визуальной обратной связи: кнопки меняют цвет при наведении, активные элементы подсвечиваются, процесс загрузки файлов отображается с помощью анимации.

\subsubsection{Организация взаимодействия с сервером}

Вся логика взаимодействия с серверной частью вынесена в отдельный модуль, реализующий единый интерфейс для выполнения HTTP-запросов. Данный модуль автоматически управляет токенами авторизации: сохраняет их в локальном хранилище браузера, добавляет в заголовки запросов и обновляет при истечении срока действия.

Для обмена данными с сервером используются следующие форматы:

\begin{enumerate}
	\item Обычные запросы (логин, регистрация, статистика) — данные в формате JSON.
	\item Загрузка файлов — данные в формате multipart/form-data.
	\item Получение сгенерированных файлов (титульных листов, проверенных документов, PDF-отчётов) — бинарные данные с последующим скачиванием.
\end{enumerate}

Взаимодействие клиентской части с сервером осуществляется посредством REST API. Для выполнения операций авторизации, загрузки документов, генерации титульных листов и получения статистики клиент отправляет HTTP-запросы к серверным маршрутам.

Основные API-запросы, используемые клиентской частью, представлены в таблице~\ref{tab:api_requests}.

\begin{xltabular}{\textwidth}{|c|X|X|X|}
	\caption{Основные REST API-запросы клиентской части\label{tab:api_requests}}\\
	\hline
	Метод & URL & Назначение & Возвращаемые данные \\ \hline
	\endfirsthead
	
	\continuecaption{Продолжение таблицы \ref{tab:api_requests}}
	\hline
	Метод & URL & Назначение & Возвращаемые данные \\ \hline
	\endhead
	
	\hline
	\endfoot
	
	POST & /auth/login & Авторизация пользователя & JWT-токены (access, refresh) \\ \hline
	POST & /auth/register & Регистрация пользователя & Данные пользователя \\ \hline
	GET & /auth/me & Получение данных текущего пользователя & ФИО, email, роль, факультет, кафедра \\ \hline
	POST & /checker/check & Загрузка и проверка документа & Результаты проверки, ссылки на отчёт и исправленный файл \\ \hline
	POST & /checker/progress/\allowbreak\{task\_id\} & Отслеживание прогресса проверки & Номер текущей страницы, общее количество, процент \\ \hline
	GET & /checker/download-checked/\{filename\} & Скачивание документа с выделенными ошибками & DOCX-файл \\ \hline
	GET & /checker/report/\allowbreak\{filename\} & Скачивание PDF-отчёта & PDF-файл \\ \hline
	POST & /title-page/generate & Генерация титульного листа методических указаний & DOCX-файл \\ \hline
	POST & /title-page/generate-tutorial & Генерация титульного листа учебного пособия & DOCX-файл \\ \hline
	POST & /title-page/generate-monograph & Генерация титульного листа монографии & DOCX-файл \\ \hline
	GET & /statistics/faculty & Статистика по факультету за период & JSON (количество документов по типам) \\ \hline
	GET & /statistics/department & Статистика по кафедре за период & JSON (количество документов по типам) \\ \hline
	GET & /statistics/user & Статистика по пользователю за период & JSON (количество документов по типам) \\ \hline
	POST & /rio/upload-step1 & Загрузка двух вспомогательных файлов (шаг 1) & Статус, переход на шаг 2 \\ \hline
	POST & /rio/upload-step2 & Загрузка и проверка основного файла (шаг 2) & Результат проверки, переход на шаг 3 \\ \hline
	POST & /rio/submit-step3 & Отправка файлов на email с комментарием (шаг 3) & Статус отправки \\ \hline
	POST & /rio/reset & Сброс сессии отправки & Статус \\ \hline
	POST & /admin/faculty & Создание факультета (только админ) & Данные созданного факультета \\ \hline
	POST & /admin/department & Создание кафедры (только админ) & Данные созданной кафедры \\ \hline
\end{xltabular}

Все сетевые вызовы выполняются асинхронно, что позволяет интерфейсу оставаться отзывчивым во время ожидания ответа сервера. При возникновении ошибок пользователю отображается понятное сообщение с пояснением причины.

\subsubsection{Реализация загрузки и проверки документов}

Ключевым элементом клиентской части является форма загрузки документов. Пользователь может переключаться между тремя типами проверяемых материалов (методические указания, учебные пособия, монографии), при этом динамически меняются текстовые подписи полей и кнопок.

Реализована поддержка перетаскивания файлов в специальную область (drag-and-drop), что повышает удобство работы. После выбора файла отображается его название, а также примерное время, необходимое для проверки (рассчитывается на основе размера файла).

При запуске проверки интерфейс переходит в состояние ожидания: кнопка отправки блокируется, отображается анимированный индикатор загрузки и прогресс-бар, показывающий текущую обрабатываемую страницу. Такой подход информирует пользователя о том, что система не зависла, а активно выполняет полезную работу.

После завершения проверки на экран выводятся результаты:

\begin{itemize}
	\item общее количество найденных ошибок;
	\item статус (ошибки найдены или документ соответствует требованиям);
	\item кнопки для скачивания проверенного документа (с выделенными ошибками) и PDF-отчёта.
\end{itemize}

\subsubsection{Генерация титульных листов}

Интерфейс генерации титульных листов поддерживает три типа документов: методические указания, учебные пособия и монографии. Для каждого типа предусмотрен свой набор полей.

Реализована возможность динамического добавления элементов формы:

\begin{enumerate}
	\item Для учебного пособия — добавление нескольких рецензентов и направлений подготовки.
	\item Для монографии — добавление нескольких авторов.
	\item Для методических указаний — добавление нескольких рецензентов.
\end{enumerate}

Каждый добавленный элемент снабжён кнопкой удаления, но система не позволяет удалить последнего рецензента или автора, чтобы гарантировать минимальную заполненность формы.

Для полей с ограничением максимальной длины (например, краткое описание) реализован счётчик символов, который динамически обновляется при вводе текста и предупреждает пользователя о превышении допустимого лимита.

\subsubsection{Отправка документов в РИО}

Для отправки документов в редакционно-издательский отдел реализован пошаговый мастер (wizard), состоящий из трёх последовательных шагов.

\textbf{Шаг 1} — загрузка двух вспомогательных файлов (произвольные документы, подтверждающие комплектность материалов).

\textbf{Шаг 2} — загрузка основного документа для проверки. В процессе загрузки система выполняет автоматическую проверку файла через серверную часть. При обнаружении ошибок пользователю выводится сообщение с требованием выполнить полную проверку файла; переход к следующему шагу блокируется.

\textbf{Шаг 3} — отображение информации об отправителе (автоматически подставляются ФИО и email текущего пользователя) и поле для ввода комментария (с ограничением 1000 символов). На этом шаге пользователь отправляет все три файла на указанный электронный адрес РИО.

Каждый шаг визуально выделяется с помощью индикатора, показывающего, какие этапы уже пройдены (зелёная галочка), какой выполняется в данный момент (синяя подсветка), а какие ещё предстоит пройти (серый цвет). Перемещение между шагами возможно только в допустимых пределах: с первого на второй — только после загрузки двух файлов, со второго на третий — только после успешного прохождения проверки.

\subsubsection{Аутентификация и управление профилем}

Для доступа к функциям системы пользователь должен пройти процедуру аутентификации. Форма входа принимает адрес электронной почты и пароль; форма регистрации дополнительно запрашивает ФИО, код факультета, название кафедры и роль.

Процесс аутентификации включает отправку учётных данных пользователя, проверку их серверной частью и получение токенов доступа. Последовательность взаимодействия компонентов представлена на рисунке~\ref{fig:auth_sequence}.

\begin{figure}[H]
	\centering
	\includegraphics[width=1\textwidth]{auth_sequence.png}
	\caption{Диаграмма последовательности авторизации пользователя}
	\label{fig:auth_sequence}
\end{figure}

После успешного входа информация о пользователе сохраняется в состоянии приложения и отображается в профиле. Пользователь может просматривать свои данные (ФИО, email, роль, факультет, кафедру, дату регистрации) и обновлять их при необходимости.

Реализован механизм выхода из системы, который отправляет на сервер запрос на отзыв токена и очищает локальное состояние приложения.

\subsubsection{Административная панель и статистика}

Пользователям с ролью администратора доступны формы для создания новых факультетов (название, код, ФИО декана) и кафедр (название, код факультета).

Интерфейс статистики позволяет формировать отчёты по трём срезам: по факультету (по коду), по кафедре (по названию) и по пользователю (по ФИО). Для каждого отчёта необходимо указать временной интервал (даты «от» и «до»). После получения ответа от сервера на экран выводятся детализированные результаты: количество методических указаний, учебных пособий, монографий и общее число документов за выбранный период.

\subsubsection{Технические особенности реализации}

Для управления состоянием приложения используется объект, хранящий текущие токены авторизации и данные профиля пользователя. Все изменения состояния автоматически синхронизируются с локальным хранилищем браузера (localStorage), что позволяет восстанавливать сеанс после перезагрузки страницы.

Реализована автоматическая обработка истечения токенов доступа: при получении от сервера ошибки авторизации клиент пытается обновить токен с помощью refresh-токена и повторяет исходный запрос. Если обновление невозможно, пользователь перенаправляется на страницу входа.

Код клиентской части организован в виде набора функций-обработчиков событий, что обеспечивает модульность и упрощает поддержку. Все функции сгруппированы по функциональным блокам: авторизация, проверка документов, генерация титульных листов, статистика, администрирование, отправка в РИО.

\subsubsection{Обеспечение безопасности}

При передаче данных на сервер используются токены авторизации, которые добавляются в заголовок Authorization каждого защищённого запроса. Access-токен хранится в локальном хранилище браузера (localStorage) и автоматически подставляется при вызове API-функций, что позволяет идентифицировать пользователя без повторной отправки пароля.

Пароли пользователей никогда не сохраняются в локальном хранилище; все операции с паролями выполняются только в зашифрованном виде при передаче на сервер, где они дополнительно хешируются с использованием алгоритма Argon2. Поля ввода пароля используют тип password, что скрывает вводимые символы от посторонних глаз.

Для защиты от подмены токенов реализован механизм их обновления: при истечении срока действия access-токена клиент автоматически отправляет refresh-токен на специальный эндпоинт, и в случае успешной проверки (токен не отозван, не истёк и принадлежит тому же пользователю) сервер выдаёт новую пару токенов. При выходе из системы оба токена удаляются из локального хранилища, а refresh-токен аннулируется на сервере (устанавливается флаг revoked).

Адрес электронной почты проходит базовую валидацию на стороне клиента перед отправкой формы, что снижает количество ошибочных запросов к серверу. Кроме того, все пользовательские данные, отображаемые в интерфейсе (например, имя файла, ФИО автора), проходят экранирование перед вставкой в HTML-код, что предотвращает атаки типа межсайтового скриптинга (XSS).

Для дополнительной защиты от несанкционированного доступа к административным функциям используются ролевые проверки: кнопка «Администрирование» и соответствующие маршруты скрыты от обычных пользователей, а любые попытки прямого вызова административных API-методов блокируются на серверной стороне с возвратом ошибки 403 Forbidden.

Все коммуникации между клиентом и сервером в рабочей среде должны осуществляться по протоколу HTTPS, что исключает перехват токенов и учётных данных при передаче по сети. 